%% kapa: exam=1 round=35 year=2024 period=2 form=A subject=appraisal_law subject_ko=감정평가관계법규 seq=2 number=2
%% source: Q-Net 감정평가사 기출문제 (cst00309) artlSeq=5213229 file=raw/1008/5213229/2교시 문제지 원본.pdf
%% conversion: kordoc 4.24.51 -> tools/gen1.py (auto); review=proofread-sonnet; flags=none
\begin{question}{2}
국토의 계획 및 이용에 관한 법령상 광역도시계획에 관한 설명으로 옳은 것은?
\begin{choices}
  \item 군수는 도지사에게 광역계획권의 지정을 요청할 수 없다.
  \item 도지사가 광역계획권을 변경하려면 중앙도시계획위원회의 심의를 거쳐 관계 중앙행정기관의 장의 승인을 받아야 한다.
  \item 국토교통부장관은 광역계획권을 변경하면 지체 없이 관계 중앙행정기관의 장에게 그 사실을 통보하여야 한다.
  \item 광역계획권을 지정한 날부터 2년이 지날 때까지 시장ㆍ군수의 광역도시계획 승인 신청이 없는 경우에는 관할 도지사가 광역도시계획을 수립한다.
  \item 국토교통부장관은 기초조사정보체계를 구축한 경우 등록된 정보의 현황을 5년마다 확인하고 변동사항을 반영하여야 한다.
\end{choices}
\end{question}
\answer{5}
